\begin{proposition}[direct boundary-Picard vanishing]
\hypertarget{tag-03IL}{}\label{wbII:prop:direct-boundary-picard}
\StatementRecord{wbII:prop:direct-boundary-picard}{proved}{03IL}{\SeriesRef{wbII:lem:picard-finite-norm}, \SeriesRef{wbII:eq:W:5.17i}, \SeriesRef{wbII:eq:W:5.17n}, \SeriesRef{wbII:eq:W:5.17o}, \SeriesRef{wbII:prop:pic0-pullback}}

\ For the standard threefold and every analytic line bundle
$L\in\operatorname{Pic}(X)$,
\[
 c_1(L|_F)=0,
 \qquad
 \eta^*i^*L\simeq\mathcal O_{\mathrm{dP}_6}.
\tag{W.41}\hypertarget{tag-03H8}{}\label{wbII:eq:W:5.17q}
\]
This proof is independent of the global Leray transgressions
$d_2^{0,q}$ and of the packaged conclusion $H^2(X;\mathbb Z)=0$; it uses
only the actual finite-filling lattices, their logarithmic overlap maps, the
based boundary relation, the local cusp solid torus, and the local integral
specialization theorem.
\end{proposition}



\begin{proof}

\begin{enumerate}[label={$\langle1\rangle$\arabic*.},leftmargin=*]
\item The common invariant calculation in \CompendiumRef{I}{wb:thm:coinvariants} first gives
\[
 E=c_1(L|_F)=kq.
\]
Because $L$ extends across the order-four filling, its fibre class belongs
to the actual specialization lattice
\[
 \pi_2^*H^2(S_2;\mathbb Z)
 =\{a\gamma\eta_2+bq:a\equiv b\pmod2\}
\]
of \CompendiumRef{I}{wb:prop:finite-indices}; \CompendiumRef{I}{wb:lem:finite-filling-stalk-map}.  Taking
coefficient independence in the displayed invariant basis forces
$a=0$ and $b=k$; hence $k$ is even.  Write $E=2nq$.

\item Contraction with the two affine translation vectors is an exact integral
calculation:
\[
 \iota_{v_1}(2q)=4\eta_1,
 \qquad
 \iota_{v_2}(2q)=-6\eta_2.
\tag{W.42}\hypertarget{tag-03H9}{}\label{wbII:eq:W:5.17r}
\]
Indeed $v_1=(1,2,-4,0)^t$, $v_2=(-1,-3,3,0)^t$ in the basis dual to
$(\gamma,u,w,\delta)$, while
$q=uw+6\gamma\delta$.  In the global coinvariants,
\[
 [\eta_1]=3[\delta],\qquad [\eta_2]=2[\delta].
\]
Applying \cref{wbII:lem:picard-finite-norm} with $m_1=3$, $m_2=4$ therefore
gives
\[
 \bar d_1=4n[\delta],\qquad
 \bar d_2=-3n[\delta].
\tag{W.43}\hypertarget{tag-03HA}{}\label{wbII:eq:W:5.17s}
\]
Now combine the balance relation \eqref{wbII:eq:W:5.17i} with the cusp vanishing \eqref{wbII:eq:W:5.17n}:
\[
 0=\bar d_1+\bar d_2+\bar d_0=n[\delta].
\]
The coinvariant group is torsion-free, so $n=0$ and hence $E=0$.

\item Finally \eqref{wbII:eq:W:5.17o} and the injectivity of the local integral specialization in
\CompendiumRef{I}{mapsI:top:thm:p0-integral-specialization} imply
$c_1(L|_W)=0$.  Pulling back to the normalization gives
$c_1(\eta^*i^*L)=0$.  The injectivity of
$c_1:\operatorname{Pic}(\mathrm{dP}_6)\to
H^2(\mathrm{dP}_6;\mathbb Z)$ proved in
\cref{wbII:prop:pic0-pullback} yields the second assertion of \eqref{wbII:eq:W:5.17q}.

\end{enumerate}
\end{proof}
